\documentclass[exercice,a4]{thedocument}%
\usepackage[french]{babel}%
\startdatakeys
    \source{21GENMATPO1}
    \BO{2026}
    \cycle{4}
    \competencesDuSocle{CA,RA}
    \objectifsApprentissage{C4C1bc,C4C2bd}
    \motsCles{probabilité,pourcentage,proportionnalité}
\enddatakeys
\begin{document}%

\medskip

Un professeur propose un jeu à ses élèves.

Ils doivent tirer un jeton dans une boîte de leur choix et gagnent lorsqu'ils tombent sur un jeton noir.

Le professeur leur précise que :

\setlength\parindent{1cm}
\begin{itemize}
\item La boîte A contient $10$ jetons dont $1$ jeton noir ;
\item La boîte B contient $15$\,\% de jetons noirs ;
\item La boîte C contient exactement $350$ jetons blancs et $50$ jetons noirs.
\end{itemize}
\setlength\parindent{0cm}

\medskip

Les jetons sont indiscernables au toucher. Une fois que l'élève a choisi sa boîte, le tirage se fait au hasard.

\medskip

\begin{enumerate}
\item Montrer que, dans la boîte C, la probabilité de tirer un jeton noir est $\dfrac{1}{8}$.
\item C'est le tour de Maxime. Dans quelle boîte a-t-il intérêt à tenter sa chance ? Justifier la réponse.
\item La boîte B contient $18$ jetons noirs. Combien y a-t-il de jetons au total dans cette boîte  ? \item On ajoute $10$ jetons noirs dans la boîte C. Déterminer le nombre de jetons blancs à ajouter dans la boîte C pour que la probabilité de tirer un jeton noir reste égale à $\dfrac{1}{8}$.
\end{enumerate}

\bigskip

\showthesource
\end{document}
